%% toc: P1 6 6.1,6.3,6.4
\section{Proofs for semantic data, models, and undischarged assumptions}
\label{sec:annex-p01-w02-p01-016-semantic-data}

Notation and definitions are those of
\cref{sec:p01-w02-p01-016-semantic-data}. Equality of assignments is
equality of functions on \(\mathsf{Var}\): two assignments are equal iff they
take the same value at every variable.

\begin{annexitem}{annex:p01-w02-p01-016-variants}{thm:p01-w02-p01-016-variants}{Updates and variants}
\begin{proof}
\emph{Clause 1.} The domain \(\dom\mathcal M\) is nonempty by the definition
of a structure; choose \(c\in\dom\mathcal M\). The function sending every
variable to \(c\) is a total function \(\mathsf{Var}\to\dom\mathcal M\),
hence an assignment.

\emph{Clause 2.} By definition \(s[x\mapsto a](y)=s(y)\) for every
\(y\ne x\), so \(s[x\mapsto a]\) is an \(x\)-variant of \(s\). Let \(s'\) be
an \(x\)-variant of \(s\). Then \(s[x\mapsto s'(x)](x)=s'(x)\), and for
\(y\ne x\), \(s[x\mapsto s'(x)](y)=s(y)=s'(y)\), the second equality because
\(s'\) is an \(x\)-variant. The two assignments agree at every variable, so
\(s'=s[x\mapsto s'(x)]\). The map \(a\mapsto s[x\mapsto a]\) therefore sends
\(\dom\mathcal M\) into the set of \(x\)-variants, and every \(x\)-variant
\(s'\) is the image of \(s'(x)\in\dom\mathcal M\), so it is onto. If
\(s[x\mapsto a]=s[x\mapsto b]\), evaluating both at \(x\) gives \(a=b\), so it
is injective.

\emph{Clause 3.} At \(x\): \(s[x\mapsto a][x\mapsto b](x)=b=s[x\mapsto b](x)\).
At \(y\ne x\): \(s[x\mapsto a][x\mapsto b](y)=s[x\mapsto a](y)=s(y)
=s[x\mapsto b](y)\). For the second identity, \(s[x\mapsto s(x)](x)=s(x)\)
and \(s[x\mapsto s(x)](y)=s(y)\) for \(y\ne x\).

\emph{Clause 4.} Let \(x\ne y\) and write \(t_1=s[x\mapsto a][y\mapsto b]\)
and \(t_2=s[y\mapsto b][x\mapsto a]\). At \(x\): since \(x\ne y\),
\(t_1(x)=s[x\mapsto a](x)=a\), and \(t_2(x)=a\). At \(y\): \(t_1(y)=b\), and
since \(y\ne x\), \(t_2(y)=s[y\mapsto b](y)=b\). At \(z\notin\{x,y\}\): both
updates leave \(z\) unchanged, so \(t_1(z)=s(z)=t_2(z)\). Hence \(t_1=t_2\).
\end{proof}
\end{annexitem}

\begin{annexitem}{annex:p01-w02-p01-016-quantifier-variants}{thm:p01-w02-p01-016-quantifier-variants}{Quantifier clauses through variants}
\begin{proof}
\emph{Clause 1, left to right.} Assume
\(\mathcal M,s\entails\forall x\,\varphi\), so
\(\mathcal M,s[x\mapsto a]\entails\varphi\) for every \(a\in\dom\mathcal M\)
by the universal satisfaction clause. Let \(s'\) be an \(x\)-variant of
\(s\). By clause 2 of \cref{thm:p01-w02-p01-016-variants},
\(s'=s[x\mapsto s'(x)]\) with \(s'(x)\in\dom\mathcal M\), so
\(\mathcal M,s'\entails\varphi\).

\emph{Clause 1, right to left.} Assume \(\mathcal M,s'\entails\varphi\) for
every \(x\)-variant \(s'\) of \(s\). For each \(a\in\dom\mathcal M\),
\(s[x\mapsto a]\) is an \(x\)-variant of \(s\) by the same clause, so
\(\mathcal M,s[x\mapsto a]\entails\varphi\). The universal satisfaction
clause gives \(\mathcal M,s\entails\forall x\,\varphi\).

\emph{Clause 2.} If \(\mathcal M,s\entails\exists x\,\varphi\), the
existential clause gives \(a\in\dom\mathcal M\) with
\(\mathcal M,s[x\mapsto a]\entails\varphi\), and \(s[x\mapsto a]\) is an
\(x\)-variant of \(s\). Conversely, if \(\mathcal M,s'\entails\varphi\) for an
\(x\)-variant \(s'\), then \(s'=s[x\mapsto s'(x)]\) and the existential clause
applies with \(a=s'(x)\).
\end{proof}
\end{annexitem}

\begin{annexitem}{annex:p01-w02-p01-016-finite-check}{thm:p01-w02-p01-016-finite-check}{Finitely many valuations suffice for propositional consequence}
\begin{proof}
For each \(u\colon P\to\mathbb B\), the function \(u^{+}\) is defined at every
member of \(\mathsf{Prop}\) and takes values in \(\mathbb B\), so it is a
valuation.

\emph{Left to right.} If \(\Gamma\entails\varphi\) and \(u^{+}\entails\Gamma\),
then \(u^{+}\entails\varphi\), because \(u^{+}\) is a valuation.

\emph{Right to left.} Assume the condition for every \(u\). Let \(v\) be a
valuation with \(v\entails\Gamma\), and put \(u=v\mathbin{\upharpoonright}P\),
the restriction of \(v\) to \(P\). Then \(u^{+}\) agrees with \(v\) on \(P\).
Every propositional variable occurring in a member \(\chi\) of
\(\Gamma\cup\{\varphi\}\) lies in \(P\), so \(u^{+}\) and \(v\) agree on the
variables of \(\chi\), and \cref{lem:prop-coincidence} gives
\(\widehat{u^{+}}(\chi)=\widehat v(\chi)\). Applied to the members of
\(\Gamma\), this gives \(u^{+}\entails\Gamma\); the assumption gives
\(u^{+}\entails\varphi\); applied to \(\varphi\), it gives
\(\widehat v(\varphi)=\widehat{u^{+}}(\varphi)=\mathbf T\). Hence
\(v\entails\varphi\), and since \(v\) was arbitrary,
\(\Gamma\entails\varphi\).

The number of functions \(P\to\mathbb B\) is \(2^{\lvert P\rvert}\), because
each of the \(\lvert P\rvert\) members of \(P\) independently receives one of
two values.
\end{proof}
\end{annexitem}

\begin{annexitem}{annex:p01-w02-p01-016-local-global}{thm:p01-w02-p01-016-local-global}{Consequence must quantify over whole tuples}
\begin{proof}
\emph{Clause 1, first statement.} Let \(\mathcal M\) satisfy
\(\mathcal M,s\entails R(x)\) for every assignment \(s\). Fix \(s\). For each
\(a\in\dom\mathcal M\), \(s[x\mapsto a]\) is an assignment, so
\(\mathcal M,s[x\mapsto a]\entails R(x)\). The universal satisfaction clause
gives \(\mathcal M,s\entails\forall x\,R(x)\).

\emph{Clause 1, second statement.} Let \(\mathcal M\) have domain \(\{0,1\}\)
and \(R^{\mathcal M}=\{0\}\), and let \(s\) send every variable to \(0\). The
value of \(x\) under \(s\) is \(0\in R^{\mathcal M}\), so
\(\mathcal M,s\entails R(x)\). The value of \(x\) under \(s[x\mapsto1]\) is
\(1\notin R^{\mathcal M}\), so \(\mathcal M,s[x\mapsto1]\not\entails R(x)\),
and the universal clause gives \(\mathcal M,s\not\entails\forall x\,R(x)\).
The datum \((\mathcal M,s)\) satisfies every member of \(\{R(x)\}\) and not
\(\forall x\,R(x)\), so \(\{R(x)\}\not\entails\forall x\,R(x)\).

\emph{Clause 2, first statement.} In this clause \(R\) is the accessibility
relation of a Kripke structure, not a relation symbol. Let
\(\mathcal M=(W,R,V)\) satisfy \(\mathcal M,w\entails p\) for every
\(w\in W\). Fix \(w\). Every \(v\) with \(wRv\) is a world, so
\(\mathcal M,v\entails p\); the clause for \(\Box\) gives
\(\mathcal M,w\entails\Box p\).

\emph{Clause 2, second statement.} Let \(W=\{0,1\}\), \(R=\{(0,1)\}\),
\(V(p)=\{0\}\) and \(V(q)=\varnothing\) for every other \(q\in\mathsf{Prop}\).
Then \(W\ne\varnothing\) and \(R\subseteq W\times W\), so \((W,R)\) is a
frame and belongs to the class of all frames. By the atomic clause,
\(\mathcal M,0\entails p\) and \(\mathcal M,1\not\entails p\). Since \(0R1\),
the clause for \(\Box\) gives \(\mathcal M,0\not\entails\Box p\). The datum
\((\mathcal M,0)\) satisfies \(p\) and not \(\Box p\), so
\(\{p\}\not\entails^{\ell}_{\mathcal C}\Box p\).
\end{proof}
\end{annexitem}

\begin{annexitem}{annex:p01-w02-p01-016-unsat-reduction}{thm:p01-w02-p01-016-unsat-reduction}{Consequence, validity and satisfiability}
\begin{proof}
\emph{The negation clause for the three semantics.} For a valuation \(v\),
\(v\entails\neg\varphi\) iff \(\widehat v(\neg\varphi)=\mathbf T\) iff
\(\widehat v(\varphi)=\mathbf F\) by the truth clause for \(\neg\), and since
\(\widehat v(\varphi)\in\{\mathbf F,\mathbf T\}\) with
\(\mathbf F\ne\mathbf T\), this holds iff
\(\widehat v(\varphi)\ne\mathbf T\), that is, iff \(v\not\entails\varphi\).
For first-order data the clause is the satisfaction clause for \(\neg\) in
\cref{sec:fol-semantics-02}. For modal data it is the classical clause for
the propositional connectives at a world.

\emph{Clause 1, left to right.} Assume \(\Gamma\entails_{\mathcal S}\varphi\)
and suppose, for contradiction, that
\(d\entails_{\mathcal S}\Gamma\cup\{\neg\varphi\}\). Then
\(d\entails_{\mathcal S}\Gamma\), so \(d\entails_{\mathcal S}\varphi\); and
\(d\entails_{\mathcal S}\neg\varphi\), so \(d\not\entails_{\mathcal S}\varphi\)
by the negation clause. This is impossible, so no datum satisfies
\(\Gamma\cup\{\neg\varphi\}\).

\emph{Clause 1, right to left.} Assume \(\Gamma\cup\{\neg\varphi\}\) is
unsatisfiable, and let \(d\entails_{\mathcal S}\Gamma\). If
\(d\not\entails_{\mathcal S}\varphi\), the negation clause gives
\(d\entails_{\mathcal S}\neg\varphi\), so \(d\) satisfies
\(\Gamma\cup\{\neg\varphi\}\), contrary to the assumption. Hence
\(d\entails_{\mathcal S}\varphi\), and \(\Gamma\entails_{\mathcal S}\varphi\).

\emph{Clause 2.} Every datum satisfies \(\varnothing\), so \(\varphi\) is
valid, that is \(\varnothing\entails_{\mathcal S}\varphi\), iff every datum
satisfies \(\varphi\). Clause 1 with \(\Gamma=\varnothing\) gives: \(\varphi\)
is valid iff \(\{\neg\varphi\}\) is unsatisfiable. For the second statement,
\(\neg\varphi\) is valid iff every datum satisfies \(\neg\varphi\), iff by the
negation clause no datum satisfies \(\varphi\), iff \(\varphi\) is
unsatisfiable. Negating both sides, \(\varphi\) is satisfiable iff
\(\neg\varphi\) is not valid.

\emph{Clause 3.} Let \(\Gamma\subseteq\Delta\). If
\(d\entails_{\mathcal S}\Delta\), then \(d\) satisfies every member of
\(\Delta\), in particular every member of \(\Gamma\), so
\(d\entails_{\mathcal S}\Gamma\). If moreover
\(\Gamma\entails_{\mathcal S}\varphi\), then \(d\entails_{\mathcal S}\varphi\);
so \(\Delta\entails_{\mathcal S}\varphi\). If \(d\) witnesses that \(\Delta\)
is satisfiable, the same \(d\) witnesses it for \(\Gamma\).

\emph{Clause 4.} If no datum satisfies \(\Gamma\), the implication
``\(d\entails_{\mathcal S}\Gamma\) implies \(d\entails_{\mathcal S}\psi\)''
holds for every datum \(d\) because its antecedent is false.

Clauses 3 and 4 used only the definitions of consequence and satisfiability,
not the negation clause.
\end{proof}
\end{annexitem}

\begin{annexitem}{annex:p01-w02-p01-016-sentences}{thm:p01-w02-p01-016-sentences}{Satisfiability and consequence for sentences through models}
\begin{proof}
Throughout we use the assignment independence corollary of
\cref{sec:fol-semantics-02}: for a sentence \(\tau\) and assignments
\(s,s'\) in \(\mathcal M\), \(\mathcal M,s\entails\tau\) iff
\(\mathcal M,s'\entails\tau\). Recall that \(\mathcal M\entails\tau\) means
\(\mathcal M,s\entails\tau\) for every assignment \(s\).

\emph{Clause 1, left to right.} Let \(\mathcal M,s\entails\tau\) for every
\(\tau\in T\). For each \(\tau\in T\) and each assignment \(s'\),
independence gives \(\mathcal M,s'\entails\tau\); so
\(\mathcal M\entails\tau\). Thus \(\mathcal M\) is a model of \(T\).

\emph{Clause 1, right to left.} Let \(\mathcal M\) be a model of \(T\). By
clause 1 of \cref{thm:p01-w02-p01-016-variants} there is an assignment \(s\)
in \(\mathcal M\). For each \(\tau\in T\), \(\mathcal M\entails\tau\) gives
\(\mathcal M,s\entails\tau\). So \((\mathcal M,s)\) satisfies \(T\).

\emph{Clause 2, left to right.} Assume \(T\entails\sigma\) and let
\(\mathcal M\in\operatorname{Mod}(T)\). For every assignment \(s\) and every
\(\tau\in T\), \(\mathcal M\entails\tau\) gives \(\mathcal M,s\entails\tau\);
so \(\mathcal M,s\entails\sigma\) by the assumption. As \(s\) was arbitrary,
\(\mathcal M\entails\sigma\), that is,
\(\mathcal M\in\operatorname{Mod}(\{\sigma\})\).

\emph{Clause 2, right to left.} Assume
\(\operatorname{Mod}(T)\subseteq\operatorname{Mod}(\{\sigma\})\), and let
\(\mathcal M,s\entails\tau\) for every \(\tau\in T\). As in clause 1,
independence gives \(\mathcal M\entails\tau\) for every \(\tau\in T\), so
\(\mathcal M\in\operatorname{Mod}(T)\), hence \(\mathcal M\entails\sigma\),
and in particular \(\mathcal M,s\entails\sigma\). So \(T\entails\sigma\).

\emph{Clause 3.} Take \(T=\varnothing\) in clause 2. Every
\(\Lang\)-structure is a model of \(\varnothing\), since the condition
``\(\mathcal M\entails\tau\) for every \(\tau\in\varnothing\)'' holds
vacuously. So \(\varnothing\entails\sigma\), which is validity of \(\sigma\),
holds iff every \(\Lang\)-structure satisfies \(\mathcal M\entails\sigma\).
\end{proof}
\end{annexitem}

\begin{annexitem}{annex:p01-w02-p01-016-closure-gap}{thm:p01-w02-p01-016-closure-gap}{Satisfiable formulas with unsatisfiable closures}
\begin{proof}
\emph{Clause 1.} Let \(\mathcal M\) have domain \(\{0,1\}\) and
\(R^{\mathcal M}=\{0\}\). Because \(x\ne y\), there is an assignment \(s\)
with \(s(x)=0\), \(s(y)=1\) and \(s(z)=0\) for every other variable \(z\).
Then \(0\in R^{\mathcal M}\) gives \(\mathcal M,s\entails R(x)\), and
\(1\notin R^{\mathcal M}\) gives \(\mathcal M,s\not\entails R(y)\), hence
\(\mathcal M,s\entails\neg R(y)\).

\emph{Clause 2.} Suppose \(\mathcal M,s\entails\forall x\,R(x)\) and
\(\mathcal M,s\entails\forall y\,\neg R(y)\). The domain is nonempty; choose
\(a\in\dom\mathcal M\). The first formula and the universal clause give
\(\mathcal M,s[x\mapsto a]\entails R(x)\); the value of \(x\) under
\(s[x\mapsto a]\) is \(a\), so \(a\in R^{\mathcal M}\). The second gives
\(\mathcal M,s[y\mapsto a]\entails\neg R(y)\), so
\(\mathcal M,s[y\mapsto a]\not\entails R(y)\) and \(a\notin R^{\mathcal M}\).
This is a contradiction, so the set is unsatisfiable. Both members are
sentences, since \(\operatorname{FV}(R(x))=\{x\}\) and the quantifier removes
\(x\), and likewise for \(y\); by clause 1 of
\cref{thm:p01-w02-p01-016-sentences}, an unsatisfiable set of sentences has
no model.

\emph{Clause 3.} Suppose \(\mathcal M\) had the property. Choose
\(a\in\dom\mathcal M\) and let \(s\) send every variable to \(a\). Then
\(\mathcal M,s\entails R(x)\) gives \(a\in R^{\mathcal M}\), and
\(\mathcal M,s\entails\neg R(y)\) gives \(a\notin R^{\mathcal M}\), a
contradiction.
\end{proof}
\end{annexitem}

\begin{annexitem}{annex:p01-w02-p01-016-open-bounds}{thm:p01-w02-p01-016-open-bounds}{How a rule application changes the open assumptions}
\begin{proof}
\emph{Clause 1.} If \(D\) is a single assumption occurrence of
\(\varphi\), that occurrence is both the leaf and the root, and the path from
it to the root contains no rule application. So it is open, and it is the
only assumption occurrence; \(\operatorname{Open}(D)=\{\varphi\}\). If
\(D\) is a single conclusion of \(=I\), it has no assumption occurrence, so
\(\operatorname{Open}(D)=\varnothing\).

\emph{Clause 2.} Since \(n\ge1\), the root of \(D\) is not a leaf, and the
leaves of \(D\) are the leaves of \(D_1,\ldots,D_n\); the subtrees are
disjoint, so each assumption occurrence \(o\) of \(D\) lies in exactly one
\(D_k\), where it is again an assumption occurrence. Whether a rule
application discharges a given occurrence depends only on the subtree whose
root that application concludes, and for every application inside \(D_k\)
this subtree is the same in \(D_k\) and in \(D\). The path from \(o\) to the
root of \(D\) is the path from \(o\) to the root of \(D_k\) followed by the
root of \(D\), which is the conclusion of \(\rho\). Hence
\[
  (\ast)\qquad o\text{ is open in }D
  \iff
  o\text{ is open in }D_k\text{ and }\rho\text{ does not discharge }o .
\]

\emph{Upper bound.} If \(\varphi\in\operatorname{Open}(D)\), some open
occurrence \(o\) of \(\varphi\) in \(D\) lies in some \(D_k\), and by
\((\ast)\) it is open in \(D_k\); so
\(\varphi\in\operatorname{Open}(D_k)\).

\emph{Lower bound.} Let \(\varphi\in\operatorname{Open}(D_k)\setminus
\Delta_\rho\), witnessed by an open occurrence \(o\) of \(\varphi\) in
\(D_k\). Because \(\varphi\notin\Delta_\rho\), \(\rho\) discharges no
occurrence of \(\varphi\), in particular not \(o\). By \((\ast)\), \(o\) is
open in \(D\), so \(\varphi\in\operatorname{Open}(D)\).

\emph{Equality.} If \(\rho\) discharges nothing, \(\Delta_\rho=\varnothing\),
and the two bounds coincide.
\end{proof}
\end{annexitem}

\begin{annexitem}{annex:p01-w02-p01-016-derivability-structure}{thm:p01-w02-p01-016-derivability-structure}{Structural properties of derivability}
\begin{proof}
\emph{Clause 1.} The one-node tree consisting of an assumption occurrence
of \(\varphi\) is a derivation with conclusion \(\varphi\), and by clause 1
of \cref{thm:p01-w02-p01-016-open-bounds} its open set is
\(\{\varphi\}\subseteq\Gamma\).

\emph{Clause 2.} A derivation \(D\) of \(\varphi\) with
\(\operatorname{Open}(D)\subseteq\Gamma\) satisfies
\(\operatorname{Open}(D)\subseteq\Delta\).

\emph{Clause 3.} Let \(D\) witness \(\Gamma\derives_{\mathcal C}\varphi\). The
tree \(D\) is finite, so it has finitely many assumption occurrences and
\(\Gamma_0\defeq\operatorname{Open}(D)\) is a finite subset of \(\Gamma\).
The same \(D\) witnesses \(\Gamma_0\derives_{\mathcal C}\varphi\).

\emph{Clause 4.} \(\derives_{\mathcal C}\varphi\) means
\(\varnothing\derives_{\mathcal C}\varphi\), that is, some derivation \(D\)
of \(\varphi\) has \(\operatorname{Open}(D)\subseteq\varnothing\). This holds
iff no formula has an open occurrence in \(D\), that is, iff every assumption
occurrence of \(D\) is discharged.
\end{proof}
\end{annexitem}

\begin{annexitem}{annex:p01-w02-p01-016-discharge}{thm:p01-w02-p01-016-discharge}{Discharge theorem for natural deduction}
\begin{proof}
Every tree built below adds rule applications only at or below the root of
an existing derivation, so each inner node of that derivation keeps its
subtree; it remains an instance of the same rule, discharges the same
occurrences, and, for \(\forall I\) and \(\exists E\) in
\(\mathsf{ND}_{\mathsf{FOL}}\), has the same open assumptions at the point of
application, so its side condition still holds. The added applications are
of \(\to E\), \(\land I\), \(\land E_2\) and \(\to I\), which have no side
conditions.

\emph{Right to left.} Let \(E\) be a derivation of \(\varphi\to\psi\) with
\(\operatorname{Open}(E)\subseteq\Gamma\). Let \(A\) be a one-node
derivation consisting of an assumption occurrence of \(\varphi\). Applying
\(\to E\) to the conclusions \(\varphi\) of \(A\) and \(\varphi\to\psi\) of
\(E\) gives a derivation \(D\) of \(\psi\). Since \(\to E\) discharges nothing,
clause 2 of \cref{thm:p01-w02-p01-016-open-bounds} holds with equality, and
clause 1 gives \(\operatorname{Open}(A)=\{\varphi\}\); so
\(\operatorname{Open}(D)=\{\varphi\}\cup\operatorname{Open}(E)
\subseteq\Gamma\cup\{\varphi\}\).

\emph{Left to right.} Let \(D\) be a derivation of \(\psi\) with
\(\operatorname{Open}(D)\subseteq\Gamma\cup\{\varphi\}\). The derivation is
finite, so some \(i\in\mathbb N\) is fresh from every label in \(D\).

\emph{Case 1: \(\varphi\in\operatorname{Open}(D)\).} Mark every open
occurrence of \(\varphi\) in \(D\) with the label \(i\). By the merging
relabeling of the inherited derivation definition \cite[ch.~``Natural
Deduction and Sequents'', definition ``Natural-deduction
derivation'']{FrozenFormalLogicSemanticData}, the result \(D'\) is a
derivation of \(\psi\) in which every rule application discharges the same
occurrences as in \(D\); hence the open occurrences of \(D'\) are those of
\(D\), and the occurrences of \(\varphi\) marked \(i\) in \(D'\) are exactly
its open occurrences of \(\varphi\), because \(i\) occurs nowhere else. Apply
\(\to I\) with label \(i\) to the conclusion \(\psi\) of \(D'\), obtaining a
derivation \(D''\) of \(\varphi\to\psi\). This application discharges exactly
the open occurrences of \(\varphi\) marked \(i\) above its premise, which are
all open occurrences of \(\varphi\) in \(D'\), and there is at least one, so
the label condition holds. By the equivalence \((\ast)\) in the proof of
\cref{thm:p01-w02-p01-016-open-bounds}, an occurrence is open in \(D''\) iff
it is open in \(D'\) and not discharged by the new root, so no occurrence of
\(\varphi\) is open in \(D''\) and every other open occurrence of \(D'\)
remains open. Thus
\(\operatorname{Open}(D'')=\operatorname{Open}(D)\setminus\{\varphi\}
\subseteq\Gamma\).

\emph{Case 2: \(\varphi\notin\operatorname{Open}(D)\).} Then
\(\operatorname{Open}(D)\subseteq\Gamma\). The rule \(\to I\) may not be
applied with a label that discharges nothing, so first introduce an occurrence
to discharge. Let \(A\) be a one-node derivation consisting of an assumption
occurrence of \(\varphi\) marked \(i\). Apply \(\land I\) to the conclusions
of \(A\) and \(D\), obtaining \(\varphi\land\psi\); apply \(\land E_2\),
obtaining \(\psi\); and apply \(\to I\) with label \(i\), obtaining
\(\varphi\to\psi\). Call the result \(D''\). The occurrence in \(A\) is open
in the subderivation ending at \(\psi\), because \(\land I\) and
\(\land E_2\) discharge nothing, and it is marked \(i\); no occurrence in
\(D\) is marked \(i\). So the final \(\to I\) discharges exactly the
occurrence in \(A\). By \((\ast)\) applied at each of the three new
applications, the open occurrences of \(D''\) are the open occurrences of
\(D\), so \(\operatorname{Open}(D'')=\operatorname{Open}(D)\subseteq\Gamma\).

In both cases \(\Gamma\derives_{\mathcal C}\varphi\to\psi\).
\end{proof}
\end{annexitem}

\begin{annexitem}{annex:p01-w02-p01-016-eigen-failure}{thm:p01-w02-p01-016-eigen-failure}{An open assumption prevents generalization}
\begin{proof}
\emph{Clause 1.} The rule \(\forall I\) may be applied to a premise only if
the generalized variable is not free in any open assumption of the
derivation of that premise \cite[ch.~``First-Order Deduction'', section
``Natural-deduction rules for quantifiers and equality'']{FrozenFormalLogicSemanticData}.
Here the derivation of the premise is the single assumption occurrence of
\(P(x)\), which is open by clause 1 of
\cref{thm:p01-w02-p01-016-open-bounds}, and
\(x\in\operatorname{FV}(P(x))\). The side condition fails, so the tree is not
a derivation. For the semantic statement, let \(\mathcal M\) have domain
\(\{0,1\}\) and \(P^{\mathcal M}=\{0\}\), and let \(s\) send every variable to
\(0\). Then \(\mathcal M,s\entails P(x)\), while
\(\mathcal M,s[x\mapsto1]\not\entails P(x)\), so
\(\mathcal M,s\not\entails\forall x\,P(x)\). Hence
\(\{P(x)\}\not\entails\forall x\,P(x)\), and by the soundness theorem
\cref{thm:fol-soundness} no derivation of \(\forall x\,P(x)\) has its open
assumptions in \(\{P(x)\}\).

\emph{Clause 2.} Take an assumption occurrence of \(P(x)\) marked \(1\) and
apply \(\to I\) with label \(1\), giving \(P(x)\to P(x)\). The occurrence is
open in the one-node subderivation and marked \(1\), so this application
discharges it, and it is the only assumption occurrence; the derivation
\(D_1\) so obtained has \(\operatorname{Open}(D_1)=\varnothing\). Apply
\(\forall I\) to the conclusion of \(D_1\). The derivation of the premise has
no open assumption, so \(x\) is free in none and the side condition holds.
The result \(D_2\) has conclusion \(\forall x\,(P(x)\to P(x))\), and since
\(\forall I\) discharges nothing, clause 2 of
\cref{thm:p01-w02-p01-016-open-bounds} gives
\(\operatorname{Open}(D_2)=\operatorname{Open}(D_1)=\varnothing\). By clause 4
of \cref{thm:p01-w02-p01-016-derivability-structure},
\(\derives_{\mathsf{ND}_{\mathsf{FOL}}}\forall x\,(P(x)\to P(x))\).
\end{proof}
\end{annexitem}
